\chapter{Estado de la Cuestión}
\label{cap:estadoDeLaCuestion}
En este capítulo se presenta una breve resumen del desarrollo de las simulaciones físicas en videojuegos, así como el hardware y software que ha hecho que esta evolución sea posible. También se presenta la simulación de telas, sus diferencias de implementación, los problemas de optimización que estas suponen y los avances en modelos de IA que permiten imaginar otras soluciones.

\section{Motores de videojuegos}
Si bien este proyecto se centra en la computación física, gestionada por los motores físicos, los motores de videojuegos proporcionan un conjunto más amplio de herramientas comunes para el desarrollo de videojuegos, entre los cuales encontramos también los motores de renderizado, sonido y otros componentes como los sistemas de scripting o interfaz de usuario, que se complementan para satisfacer las necesidades que requiere la creación de un juego. Existe gran cantidad de motores de videojuegos, siendo algunos de los más conocidos Unity y Unreal, que se detallarán a continuación. Además de estos motores comerciales de acceso general, muchos estudios grandes disponen de motores propios adaptados a sus necesidades y especialización. Algunos ejemplos son Frostbite de Electronic Arts y SnowDrop de Ubisoft. Los dos motores comerciales más utilizados son: \newline

\begin{itemize}
\item \textbf{Unreal Engine:}
Es un motor de videojuegos creado por Epic Games en 1998.  Destaca por su motor gráfico, que permite gráficos realistas en tiempo real con tecnologías como Nanite para geometría virtualizada o Lumen para iluminación dinámica. Es multiplataforma y se encuentra no sólo en pequeños desarrollos independientes (indie), sino en grandes producciones multimillonarias (producciones 'AAA'). Se caracteriza por su sistema de programación visual (Blueprints) en programación con C++. Integraba Physx como motor físico hasta la introducción reciente de Chaos Physics en su versión 5.7, que aporta novedades que se destacarán más adelante en el apartado \ref{subsec:optimtelas}.

\item \textbf{Unity:}
Creado por Unity Technologies en 2005, abarca un porcentaje del 51\% de juegos publicados en 2024 y un 26\% \footnote{Según The Big Engines Report 2025 por parte de Sensor Tower \url{https://app.sensortower.com/vgi/assets/reports/The_Big_Game_Engines_Report_of_2025.pdf}} en las unidades vendidas en Steam. Es el motor de videojuegos más utilizado por desarrolladores indie y cuenta con documentación extensa, una gran comunidad y extensiones como Sentis para la integración de modelos de IA, que explicaremos más adelante. Unity3D integra PhysX como motor de físicas, pero puede incoporar otros sistemas de físicas como Havok bajo licencia \footnote{Véasen distintas físicas integrables en Unity: \url{https://docs.unity3d.com/Manual/physics-integrations.html}}. Está basado en la arquitectura Entidad-Componente en C\#, y típicamente se utiliza para juegos más pequeños y menos potentes en cuanto a gráficos.
\end{itemize}

\section{Motores de fisica}
\label{sec:motfisicas}

Un motor de físicas es un conjunto de funciones y recursos que permite la implementación de sistemas de físicas como cuerpos rígidos con su consecuente deteccción de colisiones, cuerpos blandos e incluso fluidos. % decir algo de la estcrutra?% 
\newline

Los motores de física se extendieron en los años 90, con la popularidad del 3D y gráficos en auge. Los grandes cálculos en tiempo real con su correspondiente carga para la CPU plantearon dos grandes retos: el avance de los motores físicos cada vez más precisos y potentes, y, por otro lado, la optimización del hardware para llegar a la demanda computacional que esto suponía. \newline

En 2002 se funda Ageia Technologies, responsable de la creación del primer procesador de física dedicado para ordenadores, denominado PhysX PPU (Physics Processing Unit), así como del motor PhysX \footnote{Documentación de Physx: \url{https://developer.nvidia.com/physx-sdk}}, que sigue siendo líder en la industria a día de hoy. Posteriormente, NVIDIA Corporation, conocida por su papel en la producción de GPUs, adquiere Aegia, impulsando el cálculo de fisicas en GPU. Esto, sumado a la dificultad de penetrar en el mercado del hardware, hace que la producción de PPUs quede obsoleta. \newline % no se para que son las GPUGP)

Por otro lado, Intel adquiere en 2007 Havok \footnote{Documentación de Havok: \url{https://www.havok.com/havok-physics/}}, otro de los primeros motores de fisicas, cancelando el programa HavokFX que pretendía llevar el cálculo de fisicas a la GPU para realizar estos en sus CPUs multinúcleo. Posteriormente es adquirido por Microsoft en 2015 y se integra en las herramientas de desarollo como DirectX. Se establece como el middleware más común en grandes producciones y simulaciones complejas en estudios, mientras que Physx se convierte en libre acceso en 2018 y se integra en motores de videojuegos como Unreal.\newline

La incorporación de motores de física completos en el desarollo de videojuegos aporta escalabilidad en el \textit{gameplay}, manteniendo la coherencia de sus elementos fisicos, lo que resulta imprescindible para la inmersión del jugador \cite{hecker2000physics}.  \newline

\section{Simulación física en videojuegos}
Hasta la llegada de los motores comentados en el apartado \ref{sec:motfisicas}, existían otras técnicas para representar comportamientos físicos que eran todavía muy rudimentarias: se diseñaban las interacciones físicas previamente para parecer realistas, ya que dada la capacidad de cómputo con la que solían contar las máquinas, era imposible hacer todos los cálculos físicos en tiempo real. Se realizaban detecciones de colisiones simples, sin uso de fuerzas y descripciones de movimientos o respuestas de cada elemento del juego en lugar de sistemas de física. \newline

Estos métodos eran claramente limitados en la precisión que ofrecían, y quedaba clara la necesidad de una solución más general y reproducible. Ian Millingtong introduce en su libro \textit{Game Physics Engine Development} \cite{millington2007game} el juego \textit{Exile} \cite{exile1988}, creado por Peter Irvin y Jeremy Smith; si bien este juego fue publicado en 1988, ya contaba con conservación del momento en colisiones, gravedad, y fuerzas como el viento o explosiones. Resalta como uno de los primeros juegos que cuenta con un motor de físicas completo. \newline

Actualmente existen juegos con simulaciones físicas muy realistas como \textit{Red Dead Redemption 2} \cite{rdr2} o \textit{Hogwarts Legacy} \cite{hogwarts} donde más allá de satisfacer necesidades de gameplay, se consiguen mundos creíbles en su conjunto. Otros juegos se basan directamente en el funcionamiento y excelencia de sus físicas, particularmente simuladores de vehículos o deformaciones como la saga de \textit{Assetto Corsa} \cite{assettocorsa}, el nuevo \textit{Forza Horizon 6} \cite{forza6}, \textit{Microsoft Flight Simulator 2024} \cite{msfs2024} o \textit{BeamNG.drive} \cite{beamng}, que sigue destacando por su física basada en cuerpos blandos. \newline

Es importante conocer la diferencia entre la simulación de un cuerpo rígido y la simulación de un cuerpo blando para comprender la complejidad que implica cada uno. Se entienden los objetos deformables o cuerpos blandos como aquellos que no son rígidos, es decir, que pueden perder o modificar su forma si reciben fuerzas físicas externas. Para un cuerpo rígido, al no ser deformable, solo hay que calcular su posición, rotación y velocidad en un momento dado, en función de su estado previo y de la interacción con otros elementos físicos. Sin embargo, en el caso de un cuerpo blando, a parte de su posición y velocidad y rotación como cuerpo, hay que calcular la posición, velocidad y rotación de las partes que lo componen, a las que se les suele denominar "particulas". Estas partículas se ven afectadas por propiedades físicas como la tensión y la deformación, que modifican el comportamiento de las mismas y su interacción con otros cuerpos. Entre los cuerpos blandos se encuentran elementos como las telas, el agua o el viento \cite{erlebenbook}. \newline

Dada la complejidad de estas simulaciones, se necesita un programa que se encargue de realizar los cálculos necesarios para cada partícula. A estos programas de simulación se les denomina \textit{solvers}. Pese a encontrar múltiples acepciones en el término solver para distintas resoluciones de problemas, aquí nos referiremos como solver al algoritmo computacional que se encarga de predecir el comportamiento de estos objetos deformables en renderizados y simulaciones físicas, específicamente en la próxima sección \ref{subsec:cloth}, de telas. \newline

Otro concepto de gran importancia para este estudio es el \textit{Signed Distance Field} (SDF), que será mencionado frecuentemente en los siguientes apartados a la hora de trabajar con las colisiones con objetos externos. Los SDF representan la distancia entre un punto y la superficie más cercana de un objeto o conjunto de coordenadas. Son frecuentemente usados en la informática gráfica y generalmente se calculan como la distancia ortogonal a un espacio métrico.


\section{Simulación de telas}
El protagonista de la simulación de telas es la vestimenta: vestidos, cintas o capas moviéndose dependiendo de las acciones del jugador, viento y condiciones específicas; lo que permite a los jugadores un mayor grado de inmersión y una mayor credibilidad del mundo del videojuego. Adicionalmente, más allá de los objetos fabricados con tela, tradicionalmente ropa o banderas, las físicas de simulación de telas se pueden encontrar aplicadas en otros elementos como pelo, hierba o cuerpos blandos como podría ser un globo. Havok lo recalca en la siguiente cita: ``\textit{Havok Cloth is used to create secondary motion and to enhance any object that flows}'' \footnote{Página de Havok Cloth: \url{https://www.havok.com/havok-cloth/}}. \newline

Algunos juegos que destacan por comenzar a usar la simulación de telas en tiempo real son \textit{Assasins Creed Unity(2014)} \cite{acunity} que usó un sistema personalizado de Havok Cloth y \textit{Batman:Arkham Knight (2015)} \cite{batmanarkham} con un sistema híbrido de físicas y rigging usando Apex Cloth\footnote{Página web de Apex Clothing: \url{https://www.nvidia.com/en-us/drivers/apex-clothing/}}.

En esta sección se describe cómo se representa y simula una tela en tiempo real, el problema que supone y qué tipos de soluciones encontramos. Se indagará mayoritariamente en ejemplos de Unity3D, como referencia elegida de motor de videojuegos. 

\subsection{Cloth}
\label{subsec:cloth}
Cloth \footnote{Documentación de Cloth (Unity): \url{https://docs.unity3d.com/Manual/class-Cloth.html}}es un componente que proporciona Unity3D que funciona junto al componente \textit{Skinned Mesh Renderer} para simular telas. Está basado en Cloth de PhysX\footnote{Documentación de Cloth (Physx): \url{https://archive.docs.nvidia.com/gameworks/content/gameworkslibrary/physx/guide/3.3.4/Manual/Cloth.html}} y fue incorporado con el objetivo de poder simular el comportamiento de la ropa de personajes dentro de un videojuego. Es altamente configurable, lo que le permite adaptarse a multiples tipos de telas. Las caraterísticas parametrizables de la tela son: 
\begin{itemize}
  \item \textit{Stretching stiffness} o rigidez del estiramiento: determina cuánta resistencia opone la tela a estirarse.
  \item \textit{Bending stiffness} o rigidez de flexión: resistencia a la deformación de la tela. Un valor bajo provoca pliegues suaves y dinámicos en la tela. Junto a la rigidez de estiramiento, define la personalidad de la tela, permite recrear el material del que está formada, así como la técnica o estructura que la construye (hilos entrelazados (\textit{knit}) o continuos (tejido)).
  \item Uso de \textit{tethers}: aplica restricciones que ayudan a prevenir el movimiento de partículas de la tela de irse muy lejos de las que están fijas y reducir así el exceso de estiramiento.
  \item Uso de Gravedad sobre la tela.
  \item \textit{Damping} o coeficiente de amortiguación del movimiento: define la rapidez con la que la tela retorna a su pose de reposo.
  \item Aceleración externa y aleatoria: permiten simular fuerzas como el viento de forma integrada, con movimientos más estables con aplicando la constante de aceleración y movimientos más erráticos (como podría ser un vendaval) aplicando el valor aleatorio. 
  \item Escala de velocidad y aceleración del mundo: determinan cuánto afectará sobre los vértices de la tela la velocidad y acelaración del personaje en el espacio del mundo.
  \item Fricción: resistencia al deslizamiento que presenta la tela al colisionar con el personaje u otros colisionadores.
  \item Masa de colisión escalada: qué tanto aumentar la masa de las partículas en colisión.
  \item Uso de colisión continua: activado detecta colisiones con mayor precisión, contribuyendo a la mejora de estabilidad. 
  \item Uso de partículas virtuales: agrega una partícula virtual por triángulo para mejorar también la estabilidad de colisión.
  \item Frecuencia del solver: número de iteraciones del solver de la tela por segundo. Un número alto de iteraciones mejora las colisiones, pero sobrecarga la CPU.
  \item Umbral de adormecimiento de la tela: valor que toma el solver de la tela para considerar las partículas como quitas o ``adormecidas'' y evitar sobrecargar los cálculos.
  \item Colisionadores: dos arrays de \textit{CapsuleColliders} y \textit{SphereColliders}.
\end{itemize}

Adicionalmente, tenemos las restricciones por vértice o \textit{constraints}. Se pueden pintar a través de un mapa de calor los límites de movimiento:
\begin{itemize}
  \item \textit{Max Distance}: distancia máxima que un vértice de la tela puede moverse desde su posición original en la animación.
  \item \textit{Surface Penetration}: cuánto puede hundirse el vértice dentro de la malla del personaje.
\end{itemize}

Physx trata la tela como una malla de partículas conectadas por restricciones. El \textit{Skinned Renderer} calcula las transformaciones basándose en una malla de baja resolución sobre la que realiza los movimientos en los vértices, que dinámicamente se aplican en la malla real y se obtiene la simulación visual. Aquellos constraints no visibles o de distancia 0, fijos, no serán calculados. Cómo se ha podido apreciar en los puntos anteriores, la tela cuenta con capacidad de colisionar internamente y con dos tipos de colliders externos: cápsulas o esferas. Cualquier objeto externo o personaje con el que se quiera colisionar tendrá que ser un conjunto de estos colliders. \newline

Los primeros solvers para telas se basan en algoritmos que emplean la fuerza bruta: realizan cálculos en masa para sistemas de muelles (MSS o Mass Spring Systems) basándose en la Ley de Hooke \footnote{La ley de Hooke establece que la fuerza necesaria para deformar un cuerpo elástico, como un muelle o resorte, es directamente proporcional a la deformación producida, siempre que no se supere su límite elástico. }. Los vértices de la tela tienen un peso y están unidos por tres tipos de muelle que aseguran que la tela mantega su forma, no se distorsione o doble en exceso \cite{provot}. Los solvers más actuales resultan en el PBD o \textit{Position-Based Dynamics}, proceso iterativo que proyecta las posiciones de las partículas sobre un ``espacio de restricciones'' \cite{muller_position_2007}, es decir, actualizando el desplazamiento generado por cada fuerza en cada vertice de la simulación \cite{chaudhary2025towards}.\newline

El solver de PhysX se basa en PBD, es decir, trabaja directamente sobre las posiciones. Asímismo XPBD: Extended Position-Based \cite{macklin_xpbd_2016}, es una versión extendida que, con un coste ligeramente superior, incide en la alta dependencia sobre el paso del tiempo e iteraciones de la simulación de PBD. Unreal con Chaos integrado ya acepta constraints de tipo XPBD y también puede integrarse en Unity.


\subsection{Optimizaciones en simulación de telas}
\label{subsec:optimtelas}
Como se ha introducido antes, la simulación de telas en tiempo real supone una serie de desafíos conocidos, de los que se sigue buscando la optimización a día de hoy. El aumento de tamaño o complejidad de una tela, con su alta carga computacional, conlleva el decaímiento inmediato de frames. A parte de dificultar el realismo en juegos y demás simulaciones a tiempo real, resulta un gran impedimento en entornos dónde se requieren tasas altas de fotogramas por segundo o una menor capacidad de cómputo, por ejemplo, en dispositivos móviles sujetos a batería.\newline

Mientras que el Cloth de PhysX puede delegar la simulación a GPUs compatibles con CUDA, el Cloth integrado por defecto en Unity realiza todo su trabajo en la CPU, lo cual puede ralentizar mucho la ejecución. En el caso de Unreal, previamente al lanzamiento de ML Cloth Simulation, Chaos Physics delegaba su mayor fuerza de trabajo también en la CPU. \newline

Existen varios tipos de soluciones para la optimización de la simulación de telas:
\begin{itemize}
  \item \textbf{Optimización por CPU:} La extensión de \textit{Obi Cloth} \footnote{Obi en la Unity Asset Store: \url{https://assetstore.unity.com/packages/tools/physics/obi-cloth-81333?aid=1011l34eQ}} aprovecha el sistema de Jobs de Unity para paralelizar la simulación y acelerar la simulación sin salir de la CPU. Es uno de los paquetes más populares para solucionar este problema.
  \item \textbf{Paralelización en GPU:} a través de \textit{shaders} se delega la fuerza de cálculo del \textit{skinning} en la GPU. Diferentes estudios e individuos aportan diferentes enfoques para la optimización, por ejemplo, \cite{huguninunite16} realiza una investigación para alcanzar altas resoluciones y evitar colisiones en las mallas, mientras que otros se centran en la mejora de rendimiento pura para adaptar simulaciones a entornos VR \cite{VRGPUoptim}.
  \item \textbf{Entrenamiento con IA:}  El último enfoque es el del uso de ML para simular el comportamiento de Cloth. Esta solución ha ido cobrando fuerza en la actualidad debido a la popularización de los modelos de IA que se han aplicado con éxito en otros campos como se ha visto en la introducción \ref{cap:introduccion}. Se entrenan modelos de ML con los datos de las mallas animadas o poses, con el objetivo de que el modelo consiga simular el comportamiento de la tela utilizando menos cálculos al ejecutarlo que la simulación real.
  Se sacrifica memoria y tiempo de entrenamiento para obtener mejores y mucho más rápidas simulaciones de tela en escena. Actualmente la investigación se encuentra en el contexto de implementación para grandes empresas: Unreal ya proporciona para Chaos Cloth\footnote{Documentación introductoria de Chaos Cloth: \url{https://dev.epicgames.com/documentation/unreal-engine/machine-learning-cloth-simulation-overview}} una simulación de tela a través de aprendizaje automático, en el que destaca su realismo. Aún así, cuenta con limitaciones al basarse en un \textit{NearestNeighborSearch}: ``el sistema es capaz de predecir detalles como arrugas pero no movimiento dinámico como ondulaciones o drapeados''. En la próxima sección se detallarán las numerosas investigaciones recientes que abarcan aspectos desde la reproducción de arrugas, el compromiso entre calidad, coste y reendimiento hasta colisiones y penetraciones en las mallas.
\end{itemize}

\section{IA para optimización de telas}
Aprovechando la comprobada capacidad de las redes neuronales profundas para aproximar cualquier función matemática, una aproximación que cada vez emerge con más fuerza es el uso de estos modelos de redes neuronales para ejecutar simulaciones físicas. En estas se destacan las aplicaciones sobre telas y cuerpos no rígidos similares.

\subsection{Modelos de optimización de telas}
% Avance de simulación {Towards Generalized Position-Based Dynamics}%
Desde la introducción del PBD, la optimización de simulaciones con redes neuronales cuenta ya con un amplio recorrido. En este apartado se describen los estudios más destacados de los últimos años. Son también inspiración directa para la implementación de nuestros experimentos, con el objetivo de descubrir cuáles de estas técnicas son verdaderamente aplicables a menor escala, en desarollo indie con recursos limitados en cuanto a hardware y especialización físico-matemática.

\subsubsection{Subspace Neural Physics: Fast Data-Driven Interactive Simulation}
Este desarrollo de la unidad de investigación Ubisoft La Forge  \cite{holden_subspace_2019} implementa la primera técnica que incopora reducción de modelo o el uso de subespacios con colisiones con objetos externos en la tela. El estudio se centra en conseguir un balance entre el rendimiento en tiempo real y el coste y memoria necesarios para realizar el precómputo. \newline

El entrenamiento se realiza sobre las posiciones de una simulación realizada en Maya aplanados en un vector, al que se añade otro vector plano con los parámeteros necesarios como posición del suelo, valor del viento, o collisionadores con la tela. Se aplica un PCA\footnote{\textit{Principal Component Analysis}, es una técnica que disminuye el número de variables de un dataset, preservando la mayor cantidad de información y varianza posible y reduciendo así su complejidad} sobre ambos denotando cuántos ejes de deformación de la tela se tienen en cuenta para representar la complejidad de la recreación: 64, 128, 256; siendo este último el que conserva casi el mismo nivel de detalle que la simulación original. La arquitectura del modelo se basa en una red neuronal estándar que opera sobre el subespacio.\newline

Para evitar que la simulación colapse con el error que va produciendo, se aplican dos estretegias: 
\begin{itemize}
  \item Ventanas de frames superpuestas, con bloques de 32 frames con \textit{mini-batches}\footnote{Un \textit{batch} o lote en IA se refiere al conjunto de datos que recibe una red en su entrenamiento de forma simultánea. Su uso sirve para la optimización de recursos y memoria. } para mantener la coherencia a largo plazo
  \item Ruido gaussiano, para que la red aprenda a corregirse.
\end{itemize}
Además, los cálculos necesarios durante la ejecución se aceleran con shaders en GPU (\textit{GPU decompression}) y se recalculan las normales de los vértices necesarias para el renderizado. \newline

Sus principales limitaciones residen en la necesidad de quitar a mano errores producidos en la simulación, falta de pruebas en objetos plásticos y posibles colisiones internas fallidas. Además, se resalta como limitación realizar el entrenamiento con un objeto colisionador fijo en el espacio, sin habilitar situaciones más dinámicas como por ejemplo, un escenario en el que objetos caigan por el espacio. Finalmente se remarca la cantidad elevada de datos y tiempo empleados para los entrenamientos: se necesitaron hasta $10^{6}$ frames de simulación como referencia para obtener buenos resultados, así como entre 10 y 48 horas de entrenamiento.\newline

\subsubsection{Cloth and Skin Deformation with a Triangle Mesh Based Convolutional Neural Network}
%Our networks are built using three basic building blocks, namely Convolution, Pooling and Unpooling operators that operate directly on a manifold triangle mesh.%

Este estudio \cite{chentanez_cloth_2020} es un ejemplo de uso de CNN (redes convolucionales \footnote{Las redes neuronales convolucionales (CNN) son un tipo de arquitectura de DL diseñada para el procesamiento de datos estructurados en forma de cuadrícula, como imágenes, vídeos o secuencias temporales}), para la simulación de telas. Reseña la dificultad de intregar una CNN debido a la construcción de la malla, incompatible antes de su modificación con una cuadrícula 2D. Resuelve el problema entrenando mallas "manifold" triangulares: es decir, mallas cuyos bordes están conectados por al menos dos caras (menos los extremos) dejando una separación entre exterior e interior. Mapear la malla a una imagen 2D es lo que permite realizar la convolución. \newline

La solución planteada en este trabajo es el uso de una red encoder-decoder\footnote{\textbf{Red encoder-decoder: }Arquitectura de red neuronal diseñada para mapear una secuencia de entrada a una secuencia de salida, donde ambas pueden tener longitudes diferentes. El \textit{Encoder} comprime la información de entrada para operar con ella fácilmente, y el \textit{Decoder} procesa esa representación para generar la salida con las dimensiones reales.} formada por bloques ``DownConv'', para la reducción, y ``UpConv'', para la reconstrucción, que integra operadores nuevos de convolución\footnote{\textbf{Convolución: }Operación matemática donde un filtro o \textit{kernel} realiza multiplicaciones elemento a elemento, permitiendo extraer patrones estructurales.}. En el bloque de ``DownConv'', utiliza una técnica de pooling destacable: colapasa los edges siguiendo un método en espiral. También señala la importancia de usar Leaky-CNN: para no dejar neuronas muertas, que resulta un problema en la simulación. Además de combinar funciones de pérdida L2 y L1\footnote{\textbf{L2 y L1: }Funciones que miden la diferencia entre el valor real y el predicho. La pérdida $L_1$ (Error Absoluto Medio) calcula la suma de las diferencias absolutas, siendo robusta frente a valores atípicos. La pérdida $L_2$ (Error Cuadrático Medio) calcula la suma de los cuadrados de las diferencias, penalizando fuertemente los errores grandes.}. para que no se quede el error acumulado en ciertos vértices.\newline

Este estudio, a diferencia del anterior, incluye brazos y manos, y se rige en un marco de posiciones locales por cada triángulo.
Afirma conseguir resultados muy cercanos al \textit{ground truth}\footnote{El \textit{ground truth} es la información verdadera o correcta, a diferencia de los datos producidos por inferencia. En el caso de una recrear una simulación con ML podría consistir en el conjunto de animaciones creadas para entrenar o validar, es decir, los movimientos y comportamientos "correctos".} usando la base de animaciones \textit{Berkley animation base} y una gran capacidad de hardware disponible. No obstante, resaltan una vez más los largos tiemps de entrenamiento (entre 4 y 20 horas) y la necesidad de calidad y cantidad de datos de entrenamiento. Tampoco aporta simulaciones para topología dinámica \footnote{La topología dinámica aquí se refiere a la topología adaptativa en mallas, requeridas en aplicaciones de diseño y edición de mallas.} y se ciñe mucho al entrenamiento. Debido a la arquitectura de CNN triangular, sumado al objetivo de alcanzar velocidad y detalle en assets específicos, se pierde la capacidad de generalización.  \newline

Otro ejemplo de estudio que usa una CNN para la simulación de telas es DeepSD: Real-time Cloth Simulation with Deep Learning \cite{bertiche_deepsd_2021}; en la que proyecta la malla 3D en un plano utilizando sus coordenadas UV\footnote{Las UVs son coordenadas bidimensionales que definen cómo se proyecta una textura 2D sobre la superficie de un modelo 3D. El mapeo UV consiste en aplanar o desplegar una malla 3D sobre una superficie plana. }. Usa una CNN estándar, y no depende de una malla fija, pudiendo aceptar cambios en la topología de la malla, por lo que sirve para automatizar el proceso. El uso de SDF garantiza la generación de ropa que envuelva el cuerpo sin importar la concectividad de la malla durante los cálculos intermedios. No obstante, comparando con el estudio de  Chentanez et al. este estudio no consigue capturar arrugas finas o superficiales con detalle y es menos eficiente debido a su enfoque volumétrico. 

\subsubsection{PBNS: Physically Based Neural Simulation for Unsupervised Garment Pose Space Deformation}
PNBS \cite{bertiche2021PBNS} presenta una simulación neuronal DL no supervisada, como alternativa al enfoque clásico PBS. Integra principios físicos en el entrenamientos con funciones físicas de energía, entre ellas restricciones de fijación. La arquitectura consiste en un perceptron multicapa (MLP, por sus siglas en inglés \textit{Multi-Layer Perceptron}) con funciones ReLU\footnote{Función de activación \textit{Rectified Linear Unit} se define matemáticamente como $f(x) = \max(0, x)$} que genera deformaciones tipo PSD sobre modelos LBS, cuya aplicación en el mundo 3D garantiza la compatibilidad con los motores gráficos y su alta eficiencia. Esta metodología entrena sobre datasets públicos como \textit{AMASS} \cite{mahmood2019AMASS} y \textit{CMU MoCap}, siendo capaz de manejar colisiones en poses extremas, así como colisiones cloth-to-cloth. Asimismo, permite la obtención de  prendas que se ajustan a diferentes cuerpos (garment resizing). En comparación a otros métodos como TailorNet \cite{patel2020tailornet}, logra una mayor consistencia física y generalización, además de un coste computacional considerablemente menor, incluso sin GPU. Sin embargo, se ve limitado por su enfoque en poses estáticas.

\subsubsection{SNUG}
Siguiendo la tendencia hacia el aprendizaje no supervisado basado en físicas, la técnica SNUG \cite{santesteban2022snug} va más allá de la deformación basada en poses mencionado en el estudio anterior. Si bien trabaja con 6,519 fotogramas de la misma base de datos, gracias a una red recurrente (cuya arquitectura no se explica con exactidud) consigue mejores tiempos de entrenamiento y resultados más detallados que la técnica anterior. Pone el foco del aprendizaje en crear una serie de funciones de loss basadas en las físicas reales de deformaciones dinámicas: ``Membrane Strain Energy'' (la energía transmitida entre dos puntos en base a la elasticidad del material), ``Bending Energy'' (el ángulo diédrico entre dos caras en relación a su propiedad de rigidez), ``Gravity'' (gravedad) y ``Collision Penalty'' (la distancia prudencial entre el modelo y el cuerpo).

\subsubsection{Neural Cloth Simulation} 
Neural Cloth Simulation \cite{bertiche2022neural} también propone una metodología no supervisada, unificando simulación y entrenamiento. Sigue el esquema clásico de este tipo de solvers, donde el estado de la tela depende de los instantes previos, gestionado por ventanas de frames cuyo tamaño y complejidad dependen del tipo de tela procesada. Su arquitectura consiste en un encoder-decoder recurrente que separa componentes estáticos y dinámicos, haciendo uso de modelos GRU en estos últimos y la cual permite mejorar la generalización y estabilidad del entrenamiento. Este se realiza mediante funciones de pérdida inspirandas en simulación física, incluyendo restricciones físicas en relación a deformaciones, colisiones e inercia. El dataset utilizado también es AMASS \cite{mahmood2019AMASS}, que recopila información de mas de 40.000 poses, y el modelo físico extiende la idea del modelo mass-spring \cite{provot1995deformation}. En comparación con PBNS o SNUG, esta metodología consigue soluciones más robustas y simples que dan a simulaciones detalladas a costa de mayores tiempo de entrenamiento que sus predecesores, compensando, no obstante, la disminución de recursos al no realizar procesos de generación de datos. Esta perspectiva va de la mano con tendencias actuales en la industria, como la ya mencionada e implementada Chaos Cloth de Unreal Engine.

\subsubsection{HOOD: Hierarchical Graphs for Generalized Modelling of Clothing Dynamics}
En investigaciones más recientes, se aborda la cuestión de la posibilidad de abstraer también la topología de la tela del aprendizaje. En este caso, HOOD utiliza GNNs, una arquitectura de envío de mensajes jerárquico y por supuesto la función de pérdida basada en la física que puede verse en los previos estudios \cite{grigorev_hood_2023}.  Los graph neural networks permiten crear un gráfico que incorpora los vértices del cuerpo colisionador. A través de este se propagan los mensajes de aceleración recursivamente, lo que permite ajustar la precisión de la inferencia y adaptarse a distintos materiales, topologías y modelos de colisión. \newline 

Este estudio sigue estando limitado en aspectos como la velocidad de simulación o colisiones internas. Al fin y al cabo, este sigue siendo un campo de estudio en crecimiento, en el que todavía queda mucho por avanzar.

\subsection{Librerías }
Por último, en este apartado se describen las librerías que han resultado indispensables para entrenar los modelos de IA y conectarlos a la simulación en Unity.

\subsubsection{Pytorch}
Pytorch\footnote{Documentación de Pytorch: \url{https://docs.pytorch.org/docs/stable/index.html}} es una de las librerias de código abierto que permite implementar redes neuronales propias. Los tensores\footnote{Documentación correspondiente a los tensores: \url{https://docs.pytorch.org/docs/2.12/tensors.html}}, matrices multi-dimensionales análogas a los NumPy arrays de Python, permiten optimizar las redes tanto para el uso de CPU como de GPU . Suele ser la más elegida en entornos academicos, ya que permite prototipado rápido y entrenamiento dinámico. Fue originalmente desarrollado por Meta, y todavía cuenta con una enorme comunidad de desarrolladores.

\subsubsection{ONNX}
ONNX (\textit{Open Neural Network Exchange})\footnote{Página web de ONNX: \url{hhttps://onnx.ai/}} es un formato creado para representar modelos de aprendizaje automático . Fue introducido por Facebook y Microsoft en 2017, con el objetivo de facilitar el cambio de entornos de trabajo en aprendizaje automático. Estandariza la estructura de redes neuronales y tipos de datos, definiendo un conjunto común de operadores y un formato de archivo estándar, permitiendo entrenar desplegar modelos entrenados en un distintos ecosistemas. Es un formato ampliamente aceptado por fabricantes de hardware como NVDIA o Intel, que proporcionan proveedores de ejecución especializados para la conexión.

\subsubsection{Sentis}
Sentis (previamente Inference Engine)\footnote{Documentación del paquete Sentis: \url{https://docs.unity3d.com/Packages/com.unity.ai.inference@2.4/manual/index.html}} es una librería de inferencia de redes neuronales para Unity, predecesora introducida en el 2023 de la librería Barracuda. Permite importar modelos de redes neuronales entrenados a Unity y ejecutarlos en tiempo real tanto con la CPU como con la GPU  de forma local en el usuario, reduciendo latencia y costes de servidores en la nube. \newline

Los modelos pueden ejecutarse localmente en las plataformas que soportan Unity, lo que permite que pueda adaptarse a diferentes tipos de hardware: tanto en PC, como en móviles, consolas o web. Su compatibilidad reside en los modelos importados en formato ONNX, cuyos grafos se optimizan después para la ejecución del juego. \newline

Algunos usos de Sentis son las estimaciones de poses en tiempo real o NPCs que reaccionen a texto o voz del jugador.